\subsection*{1.2. 普通二次奇点}

设 $y$ 为特征 $p$ 的域 $k$ 上有限型概形 $Y$ 的闭点, $n$ 为 $Y$ 在 $y$ 处的维数.

\subsubsection*{定义 1.2.1 (= VI 6.6).}
\begin{enumerate}[label=(\roman*),leftmargin=2.6em]
\item \underline{若 $k$ 代数闭, 则称 $y$ 为 $Y$ 的普通二次点, 如果 $Y$ 在 $y$ 处局部环的完备化 $\widehat{\mathcal O}_{Y,y}$ 同构于 $k[[x_1,\ldots,x_{n+1}]]$ 模一个由单个形式幂级数 $f$ 生成的理想的商, 其中该幂级数可写成}
\begin{equation}\tag{1.2.1.1}
f(\underline X)=Q(\underline X)+\text{阶}>2\text{ 的项},
\end{equation}
\underline{且 $Q(\underline X)$ 为普通二次型.}
\footnote{印本中此公式编号为 (4.2.1), 且在 $Q(\underline X)$ 后多了一个等号; 英译恢复了本节编号, 并将展开式恢复为所断言的 $f=Q+$ 高次项.}

\item \underline{若 $\bar k$ 为 $k$ 的代数闭包, 则称 $y$ 为 $Y$ 的普通二次点, 如果 $\bar Y=Y\otimes_k\bar k$ 中位于 $y$ 上方的各点都是 $\bar Y$ 的普通二次点.}
\end{enumerate}

\subsubsection*{1.2.2.}
在此定义中将“普通二次型”换成“非退化二次型”, 就得到\underline{非退化二次点}的概念(VI 6.0). 点 $y$ 是 $Y$ 的非退化二次点, 当且仅当它为普通二次点, 并且 $p\ne2$ 与“$n$ 为奇数”两者至少成立其一.

\subsubsection*{例 1.2.3.}
设 $Q(\underline X)$ 为 $k$ 上的 $n+1$ 元非退化二次型. 设 $Y_0$ 为 $\mathbb A_k^{n+1}$ 中由方程 $Q=0$ 给出的二次锥. 则原点(记作 $y_0$)为 $Y_0$ 的非退化二次点.

\subsubsection*{例 1.2.4.}
假设 $p=2$ 且 $n=2m$. 设 $a\in k$, 并设 $Q$ 为多项式
\[
Q(\underline X)=(X_0^2-a)
 +\sum_{0<i\leq j\leq2m}a_{ij}x_ix_j.
\]
假设 $2m$ 元二次型 $\sum_{0<i\leq j}a_{ij}x_ix_j$ 非退化, 并以 $y_0$ 表示 $\mathbb A_k^{n+1}$ 中坐标为 $(\sqrt a,0,\ldots,0)$ 的闭点. 则 $y_0$ 为由方程 $Q=0$ 定义的闭子概形 $Y_0\subset\mathbb A_k^{n+1}$ 的普通二次点.

\subsubsection*{例 1.2.5.}
设 $k'$ 为 $k$ 的有限可分扩张, $Y_0$ 为 1.2.3(分别为 1.2.4)所考虑类型的 $k'$-概形. 若将 $Y_0$ 看作 $k$-概形, 则 $y_0$ 为 $Y_0$ 的非退化(分别为普通)二次点. 扩张 $k'/k$ 可等同于 $k(y_0)/k$ 的最大可分子扩张.

\subsubsection*{定理 1.2.6.}
\underline{假设 $y$ 为 $k$-概形 $Y$ 的普通二次点, 设 $k'$ 为 $k(y)/k$ 的最大可分子扩张. 则存在 1.2.3 或 1.2.4 中那样的 $k'$-概形 $Y_0\subset\mathbb A_{k'}^{n+1}$, 以及 $Y$ 在 $y$ 处的 Hensel 化 $Y_{(y)}$ 与 $Y_0$ 在 $y_0$ 处的 Hensel 化 $Y_{0(y_0)}$ 之间的 $k$-同构 $\varphi$.}

首先注意, $y$ 在 $Y$ 中存在一个平展邻域 $Y'$
\begin{center}
\begin{tikzcd}[row sep=large,column sep=large]
{} & Y'\arrow[d]\\
y\arrow[ur,hook]\arrow[r,hook]&Y
\end{tikzcd}
\end{center}
它是一个 $k'$-概形. 实际上, $Y_{(y)}$ 典范地是 $k'$-概形. 将 $(k,Y,y)$ 换为 $(k',Y',y)$ 后, 就归结为 $k=k'$ 的情形, 即 $k(y)/k$ 为纯不可分扩张的情形.